% GENERATED FILE. Edit canonical lessons, then run npm run generate:books.
% canonical-source-hash: fnv1a64:bf35d33a45ce517c
% canonical-lessons: ZH-C356-jiaoxun, ZH-C356-cuoguo, ZH-W356-cuo-wrong, ZH-C356-chucuo, ZH-C356-cuowu, ZH-C356-gaizheng

\chapter{A Missed Chance}
\label{ch:zh-a-missed-chance}
\hlchaptermodality{356}

\begin{chapteropening}
\textbf{By the end of this chapter:} \emph{I can say I missed a chance, slipped up and put a mistake right.}

Complete the last lesson of chapter 356: I can say I missed a chance, slipped up and put a mistake right.
\end{chapteropening}

\section[jiàoxun]{jiàoxun — hard lesson learned}

\label{lesson:ZH-C356-jiaoxun}



\begin{quote}
Before the new one: say the Mandarin for to deny, then the Mandarin for to someone's face.
\end{quote}

\subsection*{You'll want to know: jiàoxun}
\emph{jiàoxun} — \textquotedblleft{}hard lesson learned\textquotedblright{}.

\begin{quote}
\emph{wǒ juéde zhè shì yí ge hěn hǎo de jiàoxun} — I think that's a good lesson to learn.
\end{quote}

Its characters come later; for now it is a word for the ear.

\begin{cousinweb}
\emph{jiàoxun} is built from \emph{jiào} and \emph{xùn}, both inherited from Middle Chinese, \emph{kaewH} and \emph{xjunH} (Guangyun, in Baxter's transcription), which glosses the first 'to teach' and the second 'to admonish'. The second is said light here.
\end{cousinweb}

\begin{grammarlens}[title={What you've built}]
One more word for giving reasons.
\end{grammarlens}

\subsection*{Your turn}
\begin{itemize}
\raggedright
  \item \emph{Say it:} \emph{jiàoxun}
  \item \emph{Say it:} \emph{jiàoxun}, once more
  \item \emph{Say it:} \emph{wǒ juéde zhè shì yí ge hěn hǎo de jiàoxun} — I think that's a good lesson to learn.
  \item \emph{Recall:} say the Mandarin for to deny, then the Mandarin for to someone's face, then say \emph{jiàoxun} again
\end{itemize}

\begin{tcolorbox}[breakable,colback=teal!4,colframe=teal!35!black,title={Before you move on}]
What does jiàoxun mean? (\textquotedblleft{}Hard lesson learned\textquotedblright{}.) Say it once more.
\end{tcolorbox}
\section[\texorpdfstring{\zh{错过} (cuòguò)}{错过 (cuòguò)}]{\zh{错过} (cuòguò) — to miss a chance; to miss a train}

\label{lesson:ZH-C356-cuoguo}



\begin{quote}
Before the new one: say the Mandarin for to someone's face, then the Mandarin for a hard lesson learned.
\end{quote}

\subsection*{You'll want to know: \zh{错过}}
\textbf{\zh{错过}} — \emph{cuòguò} — \textquotedblleft{}to miss a chance; to miss a train\textquotedblright{}.

\begin{quote}
\textbf{\zh{我错过了那个机会}} \emph{wǒ cuòguò le nàge jīhuì} — I missed that chance.
\end{quote}

The new character, in pieces: \zh{错} — thirteen strokes in two pieces: on the left the metal piece of \zh{钱}, five strokes; on the right a short bar with two verticals through it and a long bar under them, then \zh{日} \emph{rì}, eight strokes. The next lesson puts it in your hand.

\begin{cousinweb}
\textbf{\zh{错过}} (\emph{cuòguò}) is built from \zh{错} 'wrong' and \zh{过} 'to pass', both inherited from Middle Chinese, \emph{tshak} and \emph{kwaH} (Guangyun, in Baxter's transcription), which glosses the first 'a file' and 'mixed', and the second 'a mistake' and 'to pass over'. How the first came to mean wrong is not settled.
\end{cousinweb}

\begin{grammarlens}[title={What you've built}]
One more word for giving reasons.
\end{grammarlens}

\subsection*{Your turn}
\begin{itemize}
\raggedright
  \item \emph{Say it:} \emph{cuòguò}
  \item \emph{Say it:} \emph{cuòguò}, once more
  \item \emph{Say it:} \emph{wǒ cuòguò le nàge jīhuì} — I missed that chance.
  \item \emph{Recall:} say the Mandarin for to someone's face, then the Mandarin for a hard lesson learned, then say \emph{cuòguò} again
\end{itemize}

\begin{tcolorbox}[breakable,colback=teal!4,colframe=teal!35!black,title={Before you move on}]
What does \zh{错过} mean? (\textquotedblleft{}To miss a chance; to miss a train\textquotedblright{}.) Say it once more.
\end{tcolorbox}
\section[\texorpdfstring{\zh{错} (cuò)}{错 (cuò)}]{\zh{错} — the metal piece, then a cross-bar over the sun}

\label{lesson:ZH-W356-cuo-wrong}



\begin{quote}
The word first, then the character inside it.

\begin{itemize}
\raggedright
  \item \emph{Recall:} say \emph{cuòguò} — \textbf{R1}, one lesson back, before it decays
\end{itemize}

\begin{quote}
\zh{钱} \zh{日} \zh{过}
\end{quote}

\emph{Read it:} each one aloud, with nothing said first — \emph{qián}, \emph{rì}, \emph{guò}
\end{quote}

\subsection*{Script — the metal piece, then a cross-bar over the sun}
\hlblockfigure{\detokenize{figures/ZH-W356-cuo-wrong-filmstrip.pdf}}{How \zh{错} is written, stroke by stroke}

\begin{quote}
\zh{错}
\end{quote}

\textbf{Two pieces, left before right.} Count the pieces, not the strokes.

\textbf{The left piece, the metal piece, five strokes}, as in \zh{钱}: a long left-falling stroke, a short horizontal, a second horizontal, a third that rises to the right, then a vertical that ends rising to the right.

\textbf{The right piece, eight strokes}: a short horizontal, two short verticals down through it, a long horizontal under them, then \zh{日} \emph{rì}, four strokes, at the foot.

5 + 8 = thirteen.

\subsection*{Your turn}
\begin{itemize}
\raggedright
  \item \emph{Write it:} \zh{错}, all thirteen, with the model in front of you
  \item \emph{Cover:} write it again from memory, and count the pieces, not the strokes
\end{itemize}

\begin{tcolorbox}[breakable,colback=teal!4,colframe=teal!35!black,title={Before you move on}]
How many pieces in \zh{错}? (\textbf{Two}.) Which character you know has the same left piece? (\textbf{\zh{钱}}.) What sits at the foot of the right piece? (\textbf{\zh{日}}.)

Source: \href{https://raw.githubusercontent.com/chanind/hanzi-writer-data/68d10a4b21150cae5e1ebbd223eed289cf32d90c/data/\%E9\%94\%99.json}{Hanzi Writer Data \zh{错}, pinned stroke paths and medians}.
\end{tcolorbox}
\section[\texorpdfstring{\zh{出错} (chūcuò)}{出错 (chūcuò)}]{\zh{出错} (chūcuò) — to slip up; to go wrong}

\label{lesson:ZH-C356-chucuo}



\begin{quote}
Before the new one: say the Mandarin for a hard lesson learned, then the Mandarin for to miss a chance.
\end{quote}

\subsection*{You'll want to know: \zh{出错}}
\textbf{\zh{出错}} — \emph{chūcuò} — \textquotedblleft{}to slip up; to go wrong\textquotedblright{}.

\begin{quote}
\textbf{\zh{他工作从来不出错}} \emph{tā gōngzuò cónglái bù chūcuò} — He never slips up at work.
\end{quote}

\begin{cousinweb}
\textbf{\zh{出错}} (\emph{chūcuò}) is built from \zh{出} 'to come out' and \zh{错} 'wrong', both inherited from Middle Chinese, \emph{tsyhwit} and \emph{tshak} (Guangyun, in Baxter's transcription), which glosses the first 'to go forward' and 'to appear', and the second 'mixed'.
\end{cousinweb}

\begin{grammarlens}[title={What you've built}]
One more word for giving reasons.
\end{grammarlens}

\subsection*{Your turn}
\begin{itemize}
\raggedright
  \item \emph{Say it:} \emph{chūcuò}
  \item \emph{Say it:} \emph{chūcuò}, once more
  \item \emph{Say it:} \emph{tā gōngzuò cónglái bù chūcuò} — He never slips up at work.
  \item \emph{Recall:} say the Mandarin for a hard lesson learned, then the Mandarin for to miss a chance, then say \emph{chūcuò} again
\end{itemize}

\begin{tcolorbox}[breakable,colback=teal!4,colframe=teal!35!black,title={Before you move on}]
What does \zh{出错} mean? (\textquotedblleft{}To slip up; to go wrong\textquotedblright{}.) Say it once more.
\end{tcolorbox}
\section[cuòwù]{cuòwù — mistake; error}

\label{lesson:ZH-C356-cuowu}



\begin{quote}
Before the new one: say the Mandarin for to miss a chance, then the Mandarin for to slip up.
\end{quote}

\subsection*{You'll want to know: cuòwù}
\emph{cuòwù} — \textquotedblleft{}mistake; error\textquotedblright{}.

\begin{quote}
\emph{wǒ juéde zhège cuòwù bú dà} — I don't think it's a big mistake.
\end{quote}

Its characters come later; for now it is a word for the ear.

\begin{cousinweb}
\emph{cuòwù} is built from \zh{错} 'wrong' and \emph{wù}, both inherited from Middle Chinese, \emph{tshak} and \emph{nguH} (Guangyun, in Baxter's transcription), which glosses the second 'an error'.
\end{cousinweb}

\begin{grammarlens}[title={What you've built}]
One more word for giving reasons.
\end{grammarlens}

\subsection*{Your turn}
\begin{itemize}
\raggedright
  \item \emph{Say it:} \emph{cuòwù}
  \item \emph{Say it:} \emph{cuòwù}, once more
  \item \emph{Say it:} \emph{wǒ juéde zhège cuòwù bú dà} — I don't think it's a big mistake.
  \item \emph{Recall:} say the Mandarin for to miss a chance, then the Mandarin for to slip up, then say \emph{cuòwù} again
\end{itemize}

\begin{tcolorbox}[breakable,colback=teal!4,colframe=teal!35!black,title={Before you move on}]
What does cuòwù mean? (\textquotedblleft{}Mistake; error\textquotedblright{}.) Say it once more.
\end{tcolorbox}
\section[gǎizhèng]{gǎizhèng — to put a mistake right}

\label{lesson:ZH-C356-gaizheng}



\begin{quote}
Before the new one: say the Mandarin for to slip up, then the Mandarin for an error.
\end{quote}

\subsection*{You'll want to know: gǎizhèng}
\emph{gǎizhèng} — \textquotedblleft{}to put a mistake right\textquotedblright{}.

\begin{quote}
\emph{nǐ yīnggāi gǎizhèng zhège cuòwù} — You should put this mistake right.
\end{quote}

Its characters come later; for now it is a word for the ear.

\begin{cousinweb}
\emph{gǎizhèng} is built from \emph{gǎi} and \zh{正} 'proper', both inherited from Middle Chinese, \emph{kojX} and \emph{tsyengH} (Guangyun, in Baxter's transcription), which glosses the first 'to change' and the second 'proper' and 'right'.
\end{cousinweb}

\begin{grammarlens}[title={What you've built}]
One more word for giving reasons.
\end{grammarlens}

\subsection*{Your turn}
\begin{itemize}
\raggedright
  \item \emph{Say it:} \emph{gǎizhèng}
  \item \emph{Say it:} \emph{gǎizhèng}, once more
  \item \emph{Say it:} \emph{nǐ yīnggāi gǎizhèng zhège cuòwù} — You should put this mistake right.
  \item \emph{Recall:} say the Mandarin for to slip up, then the Mandarin for an error, then say \emph{gǎizhèng} again
\end{itemize}

\begin{tcolorbox}[breakable,colback=teal!4,colframe=teal!35!black,title={Before you move on}]
What does gǎizhèng mean? (\textquotedblleft{}To put a mistake right\textquotedblright{}.) Say it once more.
\end{tcolorbox}
